\chapter{Паттерны и антипаттерны проектирования}
\label{ch:11}

Предыдущие главы дали инструменты: язык и объектную модель, коллекции,
многопоточность, сборку, базы данных, Spring, тестирование. Вы умеете написать
работающую программу. Эта глава~--- о другом: как написать программу, которую
через полгода не страшно будет менять. Разница примерно как между умением класть
кирпичи и умением проектировать дом: если строить без чертежа, через год
выяснится, что несущую стену снести нельзя, а веранду не пристроить.

\term{Паттерны проектирования} (\eng{design patterns})~--- проверенные решения
типовых задач проектирования: их не изобретают заново, их узнают и применяют.
\term{Антипаттерны} (\eng{antipatterns})~--- наоборот, решения, которые кажутся
разумными в момент написания, но систематически приводят к беде. Мы разберём
принципы SOLID, двадцать три классических паттерна и десять частых
антипаттернов~--- и выяснится, что почти все паттерны вы уже применяли: писали
\texttt{new BufferedReader(new FileReader(f))}~--- пользовались Декоратором,
ставили \texttt{@Transactional}~--- Заместителем.

\section{Что такое паттерн проектирования}

Когда один строитель говорит другому «здесь нужна несущая стена», двух слов
хватает, чтобы передать целый комплекс решений: материалы, нагрузки, технологию
монтажа. Паттерн~--- такое же слово в словаре программиста: фраза «здесь
напрашивается Стратегия» на код-ревью заменяет три абзаца объяснений о том, что
варьирующийся алгоритм надо вынести в интерфейс, а реализацию класс-клиент
должен получать извне. Формально паттерн описывается четырьмя вещами:
\term{название}~--- слово из общего словаря, \term{задача}~--- ситуация, в
которой паттерн уместен, \term{решение}~--- структура классов,
\term{последствия}~--- что вы получаете и чем за это платите. Третий пункт
важнее прочих: паттерн~--- это \emph{не готовый код для копирования}, а идея
структуры, которая в каждом проекте реализуется по-своему.

Каталог из двадцати трёх паттернов опубликовали в 1994~году Эрих Гамма, Ричард
Хелм, Ральф Джонсон и Джон Влиссидес; за авторами закрепилось прозвище
\term{«банда четырёх»} (\eng{Gang of Four, GoF}), а за паттернами~--- статус
классики. Делятся они на три группы (рис.~\ref{fig:11-1}): порождающие отвечают
на вопрос «откуда взялся объект», структурные~--- «из чего он состоит»,
поведенческие~--- «как объекты договариваются между собой».

\begin{figure}[htbp]\centering
\begin{forest} hierarchy
[Паттерны GoF
  [Порождающие
    [Singleton] [Factory Method] [Abstract Factory] [Builder] [Prototype]]
  [Структурные
    [Adapter] [Bridge] [Composite] [Decorator] [Facade] [Flyweight] [Proxy]]
  [Поведенческие
    [Chain of Responsibility] [Command] [Interpreter] [Iterator] [Mediator]
    [Memento] [Observer] [State] [Strategy] [Template Method] [Visitor]]
]
\end{forest}
\caption{Три группы паттернов «банды четырёх»}\label{fig:11-1}
\end{figure}

Частая ошибка студента, только что прочитавшего книгу по паттернам,~---
применять их везде: вместо класса на тридцать строк появляются шесть интерфейсов
и три фабрики, чтобы сложить два числа. Это \term{избыточное усложнение}
(\eng{over-engineering}): паттерн~--- лекарство, а у лекарства есть побочные
эффекты в виде лишних классов и сложной отладки. Оберегают от ошибки
\term{правило трёх} (не выносите абстракцию, пока не увидели три реальных случая
её использования) и \term{YAGNI} (\eng{You Aren't Gonna Need It}~--- не
проектируйте фабрику под пять будущих СУБД, если сейчас у вас одна). Индикатор
прост: если после введения паттерна код стало \emph{легче менять}, паттерн
нужен; если нет~--- вы усложнили на пустом месте.

\section{Принципы SOLID}

\term{SOLID}~--- акроним из пяти принципов объектно-ориентированного
проектирования, сформулированных Робертом Мартином; почти каждый паттерн GoF~---
способ соблюсти один или несколько из них. Аналогия~--- кухня ресторана: если
один повар отвечает и за салаты, и за горячее, и за мытьё посуды, кухня встанет,
как только он заболеет; если каждый отвечает за своё и общаются они через
понятные интерфейсы, кухня работает даже при замене людей.

\subsection{Единственная ответственность}

Первый принцип (\eng{Single Responsibility Principle, SRP}): у класса должна быть
только одна причина для изменения. Классический нарушитель~--- класс
\texttt{SalesReport}, который считает сумму продаж, верстает HTML-отчёт и
отправляет его почтой: правят такой класс бухгалтеры, вёрстальщики и админы
почты~--- три команды в одном файле, три источника конфликтов при слиянии.
Лечение~--- разделить по причинам изменения на \texttt{SalesCalculator},
\texttt{HtmlReportRenderer} и \texttt{ReportMailer}: строк станет больше, зато
изменение правил расчёта уже не сломает вёрстку.

\subsection{Открытость и закрытость}

Второй принцип (\eng{Open/Closed Principle, OCP}) требует, чтобы классы были
открыты для расширения, но закрыты для изменения. Аналогия~--- удлинитель с
розетками: чтобы подключить новый прибор, вы не паяете провода, а втыкаете
вилку. В листинге~\ref{lst:11-1} сравните, что придётся сделать при добавлении
треугольника в первом варианте и во втором.

\begin{listing}[H]
\begin{minted}{java}
// БЫЛО: каждая новая фигура требует правки этого метода
public double area(Object shape) {
    if (shape instanceof Circle c)
        return Math.PI * c.radius() * c.radius();
    if (shape instanceof Square s) return s.side() * s.side();
    throw new IllegalArgumentException("Неизвестная фигура");
}

// СТАЛО: поведение живёт в самих фигурах
public interface Shape { double area(); }

public record Circle(double radius) implements Shape {
    @Override public double area() { return Math.PI * radius * radius; }
}

public class AreaCalculator {
    // Не меняется никогда, сколько бы фигур ни добавили
    public double total(List<Shape> shapes) {
        return shapes.stream().mapToDouble(Shape::area).sum();
    }
}
\end{minted}
\caption{Нарушение и соблюдение принципа открытости/закрытости}
\label{lst:11-1}
\end{listing}

Добавили треугольник~--- написали новый \texttt{record Triangle implements
Shape}; существующий код не тронут, а значит, не сломан. На этом принципе
построены Стратегия, Шаблонный метод и Фабричный метод, а симптом нарушения
всегда один: длинная цепочка \texttt{instanceof} или \texttt{switch} по типу.

\subsection{Подстановка Лисков}

Третий принцип (\eng{Liskov Substitution Principle, LSP}) требует, чтобы объект
подкласса был подставим вместо объекта базового класса без нарушения работы
программы. Если у «Пингвина» метод «летать» бросает исключение, значит, пингвин
не является птицей в том смысле, в каком её понимает код, написанный для класса
«Птица». Классический пример нарушения~--- квадрат и прямоугольник
(листинг~\ref{lst:11-2}).

\begin{listing}[H]
\begin{minted}{java}
public class Rectangle {
    protected int width;
    protected int height;

    public void setWidth(int w) { this.width = w; }
    public void setHeight(int h) { this.height = h; }
    public int area() { return width * height; }
}

// С точки зрения математики квадрат - прямоугольник.
// С точки зрения кода - нет
public class Square extends Rectangle {
    @Override public void setWidth(int w) {
        this.width = w;
        this.height = w;
    }
    @Override public void setHeight(int h) {
        this.width = h;
        this.height = h;
    }
}
// r.setWidth(5); r.setHeight(4); r.area()
// для Rectangle даст 20, для Square - 16
\end{minted}
\caption{Квадрат, унаследованный от прямоугольника, нарушает контракт}
\label{lst:11-2}
\end{listing}

Заметьте: ошибка не в клиентском коде~--- он написан ровно по контракту
родителя, где ширина и высота независимы. Ошибка в наследовании, а решение~---
сделать обе фигуры неизменяемыми реализациями интерфейса \texttt{Shape}: у
неизменяемых (\eng{immutable}) объектов нет сеттеров, значит, нечему и
рассогласовываться. Наследование~--- это обещание, что ваш класс ведёт себя как
родитель; не можете его выполнить~--- используйте композицию.

\subsection{Разделение интерфейсов}

Четвёртый принцип (\eng{Interface Segregation Principle, ISP}) запрещает
заставлять клиента зависеть от методов, которые он не использует~--- это пульт с
пятьюдесятью кнопками, из которых вы жмёте четыре. Типовой нарушитель~---
«толстый» интерфейс \texttt{Device} с методами \texttt{print}, \texttt{scan} и
\texttt{fax}: простому принтеру приходится реализовывать три метода, из которых
два бросают \texttt{UnsupportedOperationException}. Лечение: одна роль~--- один
интерфейс, то есть \texttt{Printer}, \texttt{DocScanner} и \texttt{FaxSender} по
отдельности. Именно ISP объясняет обилие мелких интерфейсов в стандартной
библиотеке: \texttt{Comparable}, \texttt{Iterable}, \texttt{Runnable},
\texttt{AutoCloseable}~--- каждый описывает ровно одну роль.

\subsection{Инверсия зависимостей}

Пятый принцип (\eng{Dependency Inversion Principle, DIP}) требует, чтобы модули
верхнего уровня не зависели от модулей нижнего уровня: оба должны зависеть от
абстракций. Чайник не припаян к проводке дома~--- обе стороны зависят от
стандарта розетки, поэтому чайник можно заменить, не трогая проводку. Признак
нарушения~--- \texttt{new} с именем конкретного класса внутри бизнес-логики:
поле \texttt{MySqlOrderRepository repo = new MySqlOrderRepository()} намертво
привязывает сервис к одной базе. Лечится тем, что сервис объявляет поле типа
интерфейса \texttt{OrderRepository} и получает реализацию параметром
конструктора.

Здесь мы выходим на материал главы~\ref{ch:07}: \term{внедрение зависимостей}
(\eng{Dependency Injection}), которое делает за нас Spring,~--- техническая
реализация DIP. Сервис не знает, работает ли под ним PostgreSQL, встроенная база
в тестах или заглушка, поэтому в модульном тесте достаточно передать в
конструктор мок. Различайте три термина: \term{DIP}~--- принцип (зависеть от
абстракций), \term{DI}~--- техника (получать зависимости извне),
\term{IoC}~--- более общий принцип инверсии управления, частным случаем которого
является DI. Все пять принципов сведены в шпаргалку (табл.~\ref{tab:11-1});
правый столбец в ней важнее левого: на код-ревью вы сначала замечаете симптом и
только потом называете принцип.

\begin{table}[htbp]
\caption{Принципы SOLID и симптомы их нарушения}
\label{tab:11-1}
\centering
\begin{tabular}{L{1.6cm}L{6.4cm}L{6.8cm}}
\toprule
Принцип & Формулировка & Главный симптом нарушения \\
\midrule
SRP & одна причина для изменения & класс правят разные команды по разным поводам \\
OCP & открыт для расширения, закрыт для изменения & длинная цепочка \texttt{if}~/ \texttt{instanceof} или \texttt{switch} по типу \\
LSP & подкласс подставим вместо родителя & переопределённый метод бросает исключение или меняет смысл \\
ISP & не навязывай клиенту лишние методы & \texttt{UnsupportedOperationException} в реализации \\
DIP & зависимость от абстракций & \texttt{new ConcreteClass()} внутри бизнес-логики \\
\bottomrule
\end{tabular}
\end{table}

\section{Порождающие паттерны}

Порождающие паттерны отвечают на вопрос, как создать объект, не привязываясь
жёстко к его классу. Аналогия~--- заказ еды: можно готовить самому
(\texttt{new}), заказать блюдо в конкретной пиццерии (Фабричный метод), выбрать
целую кухню (Абстрактная фабрика) или собрать бургер по шагам в приложении
доставки (Строитель).

\subsection{Singleton (Одиночка)}

Задача~--- гарантировать, что у класса существует ровно один экземпляр, и дать к
нему глобальную точку доступа: пульт от телевизора в комнате один, и если у
каждого члена семьи заведётся свой, начнётся хаос. Наивная реализация ломается в
многопоточной среде: два потока могут одновременно пройти проверку
\texttt{if (instance == null)} и создать два объекта. Потокобезопасный вариант с
\term{двойной проверкой блокировки} (\eng{double-checked locking}) приведён в
листинге~\ref{lst:11-3}.

\begin{listing}[H]
\begin{minted}{java}
public final class ConfigRegistry {
    // volatile обязателен: без него другой поток может увидеть
    // ссылку на ещё не до конца сконструированный объект
    private static volatile ConfigRegistry instance;

    private ConfigRegistry() { /* тяжёлая инициализация */ }

    public static ConfigRegistry getInstance() {
        ConfigRegistry local = instance; // читаем поле один раз
        if (local == null) {
            synchronized (ConfigRegistry.class) {
                local = instance;
                if (local == null) {
                    instance = local = new ConfigRegistry();
                }
            }
        }
        return local;
    }
}
\end{minted}
\caption{Одиночка с двойной проверкой блокировки}
\label{lst:11-3}
\end{listing}

Слово \texttt{volatile} здесь не украшение: без него виртуальная машина вправе
переупорядочить операции так, что ссылка станет ненулевой раньше, чем
конструктор закончит работу,~--- и второй поток получит полуготовый объект
(видимость изменений между потоками разбиралась в главе~\ref{ch:05}). Проще и
надёжнее объявить Одиночку перечислением с единственной константой
\texttt{INSTANCE}: виртуальная машина сама гарантирует единственность и
потокобезопасность, а дубликат нельзя получить ни рефлексией, ни сериализацией.
Джошуа Блох называет этот способ лучшим, и в практикуме вы проверите его
устойчивость своими руками.

\begin{oshibka}
\textbf{Осторожно.} Singleton~--- самый переоценённый паттерн: глобальное
состояние затрудняет тестирование и прячет зависимости. В Spring-приложении
ручной Одиночка почти никогда не нужен: бины по умолчанию имеют область
видимости \texttt{singleton}, но живут в контейнере, а не в статическом поле,
поэтому их легко подменить в тестах.
\end{oshibka}

\subsection{Factory Method (Фабричный метод)}

Задача~--- определить интерфейс для создания объекта, но позволить подклассам
решать, какой конкретно класс создавать. Аналогия~--- сеть пиццерий: головная
компания задаёт процесс, а что именно готовит миланский филиал, решает филиал.
Абстрактный класс \texttt{NotificationDispatcher} держит общий метод
\texttt{dispatch(String text)}, вызывающий
\texttt{createNotification().send(text)}, а сам
\texttt{protected abstract Notification createNotification()} реализуют
наследники \texttt{EmailDispatcher} и \texttt{SmsDispatcher}. Настоящий
фабричный метод стандартной библиотеки~--- \texttt{Collection.iterator()}: какой
итератор вернуть, решает подкласс-коллекция. А вот
\texttt{Calendar.getInstance()}~--- не паттерн GoF, а \term{статический
фабричный метод}: выбор класса спрятан внутри одного метода, подкласс в нём не
участвует.

Рядом стоит \term{Абстрактная фабрика}, создающая целые \emph{семейства}
связанных объектов. Разница принципиальная: Фабричный метод создаёт один
продукт, Абстрактная фабрика~--- согласованный набор, как комплект мебели, где
выбранный скандинавский стиль распространяется и на стол, и на стулья. Интерфейс
\texttt{UiFactory} объявляет \texttt{createButton()} и \texttt{createCheckbox()},
а реализации \texttt{DarkUiFactory} и \texttt{LightUiFactory} возвращают
элементы своей темы; клиент знает только интерфейсы, поэтому физически не может
перемешать темы.

\subsection{Builder (Строитель)}

Задача~--- пошагово собирать сложный объект с множеством необязательных
параметров, избегая \term{телескопического конструктора}:
\texttt{new Book(title, author)}, \texttt{new Book(title, author, year)} и так
до десяти перегрузок, в которых невозможно разобраться. Аналогия~--- сборка
бургера в приложении доставки: отмечаете галочками котлету, сыр и соус и только
в конце нажимаете «Заказать». В листинге~\ref{lst:11-4} обратите внимание на две
детали: конструктор класса приватный, а вся проверка данных собрана в методе
\texttt{build}.

\begin{listing}[H]
\begin{minted}{java}
public final class Book {
    private final String title;
    private final String author;
    private final int year;
    // ... остальные поля тоже final

    private Book(Builder b) {          // конструктор приватный
        this.title = b.title;
        this.author = b.author;
        this.year = b.year;
    }

    public static Builder builder() { return new Builder(); }


    public static final class Builder {
        private String title;
        private String author;
        private int year;

        public Builder title(String t) { title = t; return this; }
        // ... author(), year() и прочие устроены так же

        public Book build() {
            // Валидация в одном месте, до создания объекта
            if (title == null || author == null) {
                throw new IllegalStateException(
                        "Название и автор обязательны");
            }
            return new Book(this);
        }
    }
}
\end{minted}
\caption{Строитель собирает неизменяемую книгу}
\label{lst:11-4}
\end{listing}

Цепочка вызовов \texttt{builder().title(...).author(...).build()} читается почти
как обычное предложение, а отдельная выгода в том, что \texttt{Book} получается
\term{неизменяемым}: поля \texttt{final}, сеттеров нет, а такие объекты
безопасны в многопоточной среде.

Наконец, \term{Прототип} создаёт новые объекты копированием существующего: если
есть заполненный бланк-образец, проще снять копию и поправить пару полей, чем
печатать бланк заново. В Java это метод \texttt{clone()} с интерфейсом
\texttt{Cloneable} либо~--- в современном коде чаще~--- \term{конструктор
копирования}. Главная ловушка~--- \term{поверхностное копирование}
(\eng{shallow copy}): \texttt{Object.clone()} копирует поля «как есть», ссылки
на вложенные объекты остаются общими, и изменение списка в копии меняет
оригинал, поэтому вложенные коллекции обязательно пересоздают строкой вида
\texttt{copy.rows = new ArrayList<>(this.rows)}.

\section{Структурные паттерны}

Структурные паттерны отвечают на вопрос, как собрать из объектов более крупную
конструкцию, сохранив гибкость. Аналогия~--- конструктор: одни детали соединяют
несовместимое, другие оборачивают деталь, добавляя ей свойства, третьи собирают
деревья из одинаковых элементов.

\subsection{Adapter и Bridge (Адаптер и Мост)}

Задача \term{Адаптера}~--- заставить работать вместе классы с несовместимыми
интерфейсами. Аналогия прямо в названии: вилка европейская, розетка британская,
менять ни то ни другое вы не можете~--- берёте переходник. Устройство всегда
одно: класс \texttt{XmlLoggerAdapter} реализует нужный приложению интерфейс
\texttt{AppLogger}, хранит в поле объект чужой библиотеки и в методе \texttt{log}
переводит вызов на её язык. В стандартной библиотеке адаптеры~---
\texttt{InputStreamReader} (байтовый поток в символьный, глава~\ref{ch:05}) и
\texttt{Arrays.asList()}.

\term{Мост} разделяет абстракцию и реализацию, решая проблему комбинаторного
взрыва наследования: пусть есть обычные и срочные уведомления, а отправлять их
можно по почте и в мессенджер~--- через наследование получится четыре класса;
добавили SMS~--- шесть; добавили отложенные уведомления~--- девять. Технически
мост~--- это поле: абстрактный класс \texttt{Notification} хранит ссылку
\texttt{protected final MessageSender sender} и получает её в конструкторе,
наследники вызывают \texttt{sender.send(...)}, ничего не зная о канале, а
реализации \texttt{EmailSender} и \texttt{TelegramSender}~--- о видах
уведомлений. Произведение превращается в сумму: три вида плюс три канала~---
шесть классов вместо девяти. Самый известный мост в мире Java~--- архитектура
JDBC: \texttt{Connection} и \texttt{Statement} являются абстракцией, а драйвер
конкретной СУБД~--- реализацией (глава~\ref{ch:06}).

\subsection{Composite (Компоновщик)}

Задача~--- работать с деревом объектов так же, как с отдельным объектом.
Аналогия~--- файловая система: папка содержит файлы и другие папки, но когда вы
спрашиваете «сколько это весит», вам всё равно, папка перед вами или файл. Ключ
в том, что и лист дерева, и составной узел реализуют один интерфейс: объявите
\texttt{interface FileNode \{ String name(); long size(); \}}, сделайте лист
записью \texttt{record TextFile(String name, long size) implements FileNode}, а
составной узел \texttt{Folder} пусть хранит \texttt{List<FileNode> children} и в
своём \texttt{size()} суммирует \texttt{child.size()} по всем детям, рекурсивно
спускаясь по дереву. Так устроен граф сцены в JavaFX (\texttt{Node} и
\texttt{Parent}): контейнер сам является узлом, поэтому вложенность
произвольная.

\subsection{Decorator (Декоратор)}

Задача~--- динамически добавлять объекту новые обязанности, не меняя его класс и
не плодя подклассы. Аналогия~--- кофе: есть эспрессо, добавили молоко~--- латте,
добавили сироп~--- латте с сиропом; каждая добавка оборачивает напиток, но
напиток остаётся напитком. В листинге~\ref{lst:11-6} декораторы реализуют тот же
интерфейс, что и оборачиваемый объект, и хранят ссылку на него.

\begin{listing}[H]
\begin{minted}{java}
public interface DataSource { String read(); }

public class FileDataSource implements DataSource {
    @Override public String read() { return "данные отчёта"; }
}

// Базовый декоратор: тот же интерфейс плюс вложенный объект
public abstract class DataSourceDecorator implements DataSource {
    protected final DataSource wrappee;
    protected DataSourceDecorator(DataSource w) { wrappee = w; }
}

public class UpperCaseDecorator extends DataSourceDecorator {
    public UpperCaseDecorator(DataSource w) { super(w); }
    @Override public String read() {
        return wrappee.read().toUpperCase();
    }
}

// BracketsDecorator устроен так же, но его read()
// возвращает "[" + wrappee.read() + "]"
// new BracketsDecorator(new UpperCaseDecorator(
//         new FileDataSource())).read()
// вернёт [ДАННЫЕ ОТЧЁТА]
\end{minted}
\caption{Декораторы складываются как матрёшка}
\label{lst:11-6}
\end{listing}

Это самый узнаваемый паттерн стандартной библиотеки: \texttt{BufferedReader}
добавляет буферизацию и метод \texttt{readLine()}, оставаясь потомком
\texttt{Reader}; так же устроен \texttt{Collections.unmodifiableList()}~---
только он не добавляет возможность, а отнимает. Отличие от наследования в том,
что декораторы собираются во время работы программы в любом порядке, тогда как
наследованием пришлось бы завести по классу на каждое сочетание свойств.

\subsection{Facade и Flyweight (Фасад и Приспособленец)}

Задача \term{Фасада}~--- предоставить простой интерфейс к сложной подсистеме,
как оператор колл-центра банка, за которым стоят десятки систем (скоринг,
процессинг, антифрод), а вы говорите «хочу оформить карту» и не знаете о них
ничего. Класс \texttt{OrderFacade} с единственным методом \texttt{placeOrder}
вызывает три подсистемы в правильном порядке: резервирует книгу, списывает
деньги, назначает доставку. Соблазн написать \texttt{new InventoryService()}
прямо в поле велик, но это тот самый симптом нарушения DIP: подсистемы должны
приходить в конструкторе, тогда фасад можно протестировать с заглушками. И
важно: фасад не запрещает обращаться к подсистемам напрямую, он лишь предлагает
удобный путь для типового сценария. В Spring фасад~--- \texttt{JdbcTemplate}.

\term{Приспособленец} экономит память, разделяя общее неизменяемое состояние
между множеством объектов, как типографская касса букв: чтобы набрать книгу, не
отливают отдельную литеру для каждой буквы, берут одну и ставят в нужные места;
внешнее состояние (позиция на странице) хранится отдельно, внутреннее (форма
буквы)~--- общее. Реализация сводится к фабрике с кешем: метод
\texttt{CharStyle of(String font, int size, String color)} собирает ключ из трёх
параметров и возвращает объект из общей карты через \texttt{computeIfAbsent}, так
что на документе в миллион символов с пятью стилями в памяти будет пять
объектов. В стандартной библиотеке приспособленец объясняет знаменитую загадку:
\texttt{Integer.valueOf()} кеширует значения от $-128$ до $127$, поэтому для
\texttt{Integer a = 127, b = 127} сравнение \texttt{a == b} даёт \texttt{true},
а для 128~--- уже \texttt{false}. Так же работает пул строковых литералов.

\subsection{Proxy (Заместитель)}

Задача~--- подставить вместо объекта его дублёра с тем же интерфейсом, чтобы
контролировать доступ к оригиналу: секретарь руководителя для посетителя говорит
от лица начальника, часть вопросов решает сам, а до начальника доводит только
нужное. В листинге~\ref{lst:11-7} заместитель кеширует результаты, и оригинал об
этом ничего не знает.

\begin{listing}[H]
\begin{minted}{java}
public interface BookRepository { String findTitle(long id); }

public class RealBookRepository implements BookRepository {
    @Override public String findTitle(long id) {
        System.out.println("Тяжёлый запрос в БД, id=" + id);
        return "Книга №" + id;
    }
}

// Тот же интерфейс, дополнительное поведение
public class CachingBookProxy implements BookRepository {
    private final BookRepository target;
    private final Map<Long, String> cache = new HashMap<>();

    public CachingBookProxy(BookRepository t) { target = t; }

    @Override public String findTitle(long id) {
        return cache.computeIfAbsent(id, target::findTitle);
    }
}
\end{minted}
\caption{Кеширующий заместитель репозитория}
\label{lst:11-7}
\end{listing}

Разновидностей заместителя несколько: кеширующий, защитный (проверяет права),
ленивый (создаёт тяжёлый объект при первом обращении), удалённый (прячет сетевой
вызов) и логирующий. Технически Заместитель почти не отличается от
Декоратора~--- оба реализуют интерфейс и держат ссылку на объект,~--- но
различаются намерением: Декоратор \emph{добавляет возможности} осознанно для
клиента, а Заместитель \emph{контролирует доступ}, и клиент часто даже не знает
о его существовании.

Это ключевой паттерн всего Spring: класс \texttt{java.lang.reflect.Proxy} создаёт
заместителя по интерфейсу прямо во время работы программы, и на этом механизме
построены \texttt{@Transactional}, \texttt{@Cacheable} и вся аспектная механика.
Отсюда следует правило, которое в главе~\ref{ch:07} приходилось принимать на
веру: \texttt{@Transactional} не срабатывает при вызове метода изнутри того же
класса, потому что вызов \texttt{this.method()} идёт мимо заместителя
(рис.~\ref{fig:11-2}).

\begin{figure}[htbp]\centering
\begin{tikzpicture}[font=\small,node distance=8mm]
\node[boxw=20mm] (cl) {клиент\\[1pt]\scriptsize (другой бин)};
\node[boxw=32mm,right=of cl] (px)
  {\textbf{Заместитель}\\[1pt]\scriptsize открыть транзакцию,\\
   \scriptsize вызвать, закоммитить};
\node[boxw=27mm,right=of px] (tg)
  {\texttt{BookService}\\[1pt]\scriptsize ваш настоящий бин};
\draw[flowarrow] (cl)--(px);
\draw[flowarrow] (px)--(tg);
\draw[flowarrow] (tg.south) .. controls +(0,-9mm) and +(9mm,-9mm) .. (tg.east);
\node[font=\scriptsize,below=8mm of tg.south,anchor=north,align=center]
  {\texttt{this.method()} идёт мимо заместителя};
\end{tikzpicture}
\caption{Заместитель Spring и внутренний вызов метода}\label{fig:11-2}
\end{figure}

\section{Поведенческие паттерны}

Поведенческие паттерны~--- про распределение обязанностей и способы общения
объектов. Аналогия~--- организация работы в компании: кто кому подчиняется, как
передаётся заявка по инстанциям, кто кого уведомляет об изменениях, кто решает,
каким способом выполнять задачу.

\subsection{Шесть паттернов, которые вы уже применяли}

Эти шесть паттернов не требуют отдельных листингов: вы либо уже пользовались
ими, либо встретите их в готовом виде.

\term{Цепочка обязанностей} передаёт запрос по цепочке обработчиков, пока один
из них его не обработает: заявление в вузе попадает в деканат, оттуда в учебную
часть, оттуда к проректору, и каждая инстанция либо решает вопрос сама, либо
передаёт дальше. Базовый класс \texttt{Validator} хранит ссылку на следующее
звено и в финальном методе \texttt{validate} сначала вызывает свою проверку, а
затем, если она прошла, обращается к следующему звену; метод \texttt{linkWith}
обычно возвращает присоединённое звено, поэтому в переменной нужно держать
первое~--- именно ему отправляют запрос. Так устроены цепочка фильтров сервлетов
и \texttt{SecurityFilterChain} в Spring.

\term{Команда} превращает запрос в объект, чтобы его можно было передать,
поставить в очередь, залогировать или отменить,~--- как заказ, записанный
официантом на бумажке. Интерфейс \texttt{Command} с методами \texttt{execute()}
и \texttt{undo()} плюс класс \texttt{CommandHistory} со стеком выполненных
команд~--- вот и весь паттерн; \texttt{Runnable} и \texttt{Callable}~--- это и
есть команды, которые вы кладёте в пул потоков методом \texttt{submit()}
(глава~\ref{ch:05}).

\term{Итератор} даёт последовательный доступ к элементам составного объекта, не
раскрывая его устройство: вы не изучаете план музея, а идёте за экскурсоводом.
Весь фреймворк коллекций построен на интерфейсах \texttt{Iterator} и
\texttt{Iterable}, а цикл \texttt{for-each}~--- сахар над итератором.
\term{Посредник} убирает связи «каждый с каждым», заставляя объекты общаться
через единый объект: самолёты не договариваются о посадке напрямую, все общаются
с диспетчерской вышкой; так работают \texttt{ExecutorService} и
\texttt{ApplicationEventPublisher}. \term{Снимок} (\eng{Memento}) сохраняет
состояние объекта, чтобы позже его восстановить, не нарушая инкапсуляцию,~---
как сохранение в игре перед сложным боем; в Java снимок удобно делать записью и
хранить в стеке. Наконец, \term{Интерпретатор} задаёт грамматику языка и
вычисляет выражения, представляя их деревом объектов: так устроены
\texttt{java.util.regex.Pattern} и язык выражений SpEL, но применим он только к
простым грамматикам~--- для настоящего языка пишут парсер генератором.

\subsection{Observer (Наблюдатель)}

Задача~--- оповещать множество объектов об изменении состояния другого объекта,
не связывая их жёстко: автор канала публикует ролик и не знает поимённо, кто его
посмотрит, а подписчики получают уведомление автоматически и могут отписаться в
любой момент. Листинг~\ref{lst:11-8} показывает издателя целиком.

\begin{listing}[H]
\begin{minted}{java}
@FunctionalInterface
public interface OrderListener {
    void onOrderCreated(String orderId);
}

public class OrderService {
    private final List<OrderListener> listeners =
            new ArrayList<>();

    public void subscribe(OrderListener l) { listeners.add(l); }
    public void unsubscribe(OrderListener l) {
        listeners.remove(l);
    }

    public void createOrder(String orderId) {
        System.out.println("Заказ " + orderId + " создан");
        for (OrderListener listener : listeners) {
            listener.onOrderCreated(orderId);
        }
    }
}
// service.subscribe(id -> System.out.println("Письмо " + id));
\end{minted}
\caption{Издатель рассылает события подписчикам}
\label{lst:11-8}
\end{listing}

Ключевая выгода: чтобы добавить новую реакцию на событие, класс
\texttt{OrderService} менять не нужно~--- это принцип OCP в действии.

\begin{oshibka}
\textbf{Две классические ошибки Наблюдателя.} Первая: слушатель отписывается
прямо внутри обработчика~--- список меняется во время перебора, и вы получаете
\texttt{ConcurrentModificationException}; лечится перебором копии списка или
хранением слушателей в \texttt{CopyOnWriteArrayList}. Вторая: один упавший
слушатель обрывает рассылку остальным~--- лечится обёрткой каждого вызова в
\texttt{try/catch}. Обе ошибки вы воспроизведёте в практикуме.
\end{oshibka}

\subsection{State (Состояние)}

Задача~--- позволить объекту менять поведение при изменении внутреннего
состояния, так, будто он меняет класс: турникет в метро на толчок отвечает
«заблокировано», а после оплаты тот же толчок пропускает пассажира. В
листинге~\ref{lst:11-9} каждое состояние заказа~--- отдельный класс, знающий,
куда из него можно перейти.

\begin{listing}[H]
\begin{minted}{java}
public interface OrderState {
    OrderState pay();
    OrderState cancel();
    String title();
}

public class NewState implements OrderState {
    @Override public OrderState pay() { return new PaidState(); }
    @Override public OrderState cancel() {
        return new CancelledState();
    }
    @Override public String title() { return "НОВЫЙ"; }
}

public class PaidState implements OrderState {
    // Из оплаченного платить второй раз нельзя,
    // а отменить - только с возвратом средств
    @Override public OrderState pay() {
        throw new IllegalStateException("Заказ уже оплачен");
    }
    // ... cancel() возвращает new CancelledState()
    @Override public String title() { return "ОПЛАЧЕН"; }
}

public class Order {
    private OrderState state = new NewState();

    public void pay() { state = state.pay(); }
    public void cancel() { state = state.cancel(); }
    public String status() { return state.title(); }
}
\end{minted}
\caption{Каждое состояние заказа - отдельный класс}
\label{lst:11-9}
\end{listing}

Сравните с альтернативой~--- полем \texttt{int status} и цепочкой
\texttt{if (status == 1) ... else if (status == 2)}: такой код растёт при каждом
новом статусе и склонен к ошибкам, ведь статусов может быть уже семь, а условие
проверяет пять.

\subsection{Strategy (Стратегия)}

Задача~--- вынести семейство взаимозаменяемых алгоритмов в отдельные объекты и
позволить подменять их на лету: навигатор строит маршрут до одной точки
по-разному («самый быстрый», «без платных дорог», «пешком»). Благодаря
функциональным интерфейсам из главы~\ref{ch:06} стратегия пишется одной строкой
(листинг~\ref{lst:11-10}).

\begin{listing}[H]
\begin{minted}{java}
@FunctionalInterface
public interface DiscountStrategy { double apply(double sum); }

public class Cart {
    // Стратегия по умолчанию: без скидки
    private DiscountStrategy strategy = sum -> sum;

    public void setStrategy(DiscountStrategy s) { strategy = s; }

    public double total(double sum) {
        return strategy.apply(sum);
    }
}
// cart.total(1000)                     вернёт 1000.0
// cart.setStrategy(s -> s * 0.9);  ->  total(1000) = 900.0
\end{minted}
\caption{Стратегия расчёта скидки}
\label{lst:11-10}
\end{listing}

Структурно Стратегия почти не отличается от Состояния, но различается намерением:
в Стратегии алгоритм выбирает \emph{клиент} извне, и стратегии не знают друг о
друге; в Состоянии следующий переход выбирает \emph{сам объект состояния}.
Классическая стратегия стандартной библиотеки~--- \texttt{Comparator}; более
того, любой бин, внедряемый через интерфейс,~--- это Стратегия, а контейнер
выступает механизмом её подстановки.

\subsection{Template Method (Шаблонный метод)}

Задача~--- задать скелет алгоритма в базовом классе, позволив подклассам
переопределять отдельные шаги: в рецепте горячего напитка шаги всегда одни
(вскипятить воду, заварить, налить в чашку), но «заварить» для чая и кофе
означает разное. Абстрактный класс \texttt{DataImporter} объявляет
\texttt{public final void importData(String source)}, где вызывает
\texttt{open(source)}, затем в блоке \texttt{try} перебирает строки из
\texttt{readRows()} методом \texttt{process(row)}, а в \texttt{finally}
закрывает источник. Шаги \texttt{open}, \texttt{readRows} и \texttt{process}
объявлены \texttt{protected abstract} и реализуются в наследниках, а
\texttt{close()}~--- \term{хук} (\eng{hook}) с пустой реализацией по умолчанию.
Слово \texttt{final} на самом шаблонном методе~--- не формальность: оно защищает
инвариант, оставляя порядок шагов и обработку ошибок за базовым классом. Так
устроен \texttt{HttpServlet.service()}, вызывающий \texttt{doGet()} или
\texttt{doPost()} в зависимости от метода запроса, а в Spring~---
\texttt{JdbcTemplate} и \texttt{TransactionTemplate}.

\subsection{Visitor (Посетитель)}

Задача~--- добавить новую операцию над объектами структуры, не меняя классы этих
объектов; аналогия~--- ревизор, обходящий отделы компании: отделы не переписывают
регламенты под каждую проверку, они просто принимают ревизора. В классическом
виде паттерн требует метода \texttt{accept(Visitor)} в каждом узле и метода
\texttt{visit} на каждый тип узла в посетителе; цена такова: новую операцию
добавить легко, а новый тип узла тяжело. В современной Java есть более короткая
альтернатива~--- запечатанные (\texttt{sealed}) интерфейсы из главы~\ref{ch:03}
вместе с сопоставлением с образцом в \texttt{switch}
(листинг~\ref{lst:11-12}).

\begin{listing}[H]
\begin{minted}{java}
public sealed interface Node permits Text, Image { }

public record Text(String value) implements Node { }
public record Image(String url) implements Node { }

public String render(Node node) {
    return switch (node) {          // ветка default не нужна
        case Text t -> "<p>" + t.value() + "</p>";
        case Image i -> "<img src=\"" + i.url() + "\">";
    };
}
\end{minted}
\caption{Обход структуры через sealed-интерфейс и switch}
\label{lst:11-12}
\end{listing}

Компилятор сам проверит, что обработаны все варианты: слово \texttt{sealed}
гарантирует, что других наследников нет.

\section{Паттерны вокруг нас: стандартная библиотека и Spring}

Пока вы не знали имён, стандартная библиотека была ровным фоном, а стоит выучить
имена~--- и в каждой второй строке проступает знакомый силуэт. Возьмём цепочку
потоков ввода-вывода из главы~\ref{ch:05}:
\texttt{new BufferedReader(new InputStreamReader(new
FileInputStream("data.txt"), UTF\_8))}. Здесь два разных паттерна в одном
выражении: \texttt{FileInputStream}~--- источник байтов, паттерна за ним нет;
\texttt{InputStreamReader}~--- \emph{Адаптер}, превращающий байтовый поток в
символьный (тип на входе и на выходе разный); \texttt{BufferedReader}~---
\emph{Декоратор}, добавляющий буферизацию и метод \texttt{readLine()} и
остающийся при этом \texttt{Reader}. Это и есть признак различения: адаптер
меняет интерфейс, декоратор сохраняет. А в вызове
\texttt{jdbcTemplate.query(sql, (rs, n) -> rs.getString("title"), 2000)}
паттернов сразу три: \texttt{JdbcTemplate}~--- одновременно Шаблонный метод
(неизменную часть работы с базой он делает сам) и Фасад над JDBC API, а
переданная лямбда~--- Стратегия отображения строки результата в объект.

Сводка «какой паттерн где искать» приведена в табл.~\ref{tab:11-2}: если по
строке таблицы вы можете за минуту вспомнить структуру паттерна, материал
усвоен.

\begin{longtable}{L{3.4cm}L{11.6cm}}
\caption{Паттерны в стандартной библиотеке и в Spring}
\label{tab:11-2}\\
\toprule
Паттерн & Где встречается \\
\midrule
\endfirsthead
\toprule
Паттерн & Где встречается \\
\midrule
\endhead
\bottomrule
\endfoot
Singleton & \texttt{Runtime.getRuntime()}; область видимости \texttt{singleton} у бинов Spring \\
Factory Method & \texttt{Collection.iterator()}; интерфейс \texttt{FactoryBean} \\
Abstract Factory & \texttt{DocumentBuilderFactory}, \texttt{TransformerFactory} \\
Builder & \texttt{StringBuilder}, \texttt{HttpRequest.newBuilder()}, \texttt{HttpSecurity} \\
Prototype & \texttt{ArrayList.clone()}; область видимости \texttt{prototype} \\
Adapter & \texttt{InputStreamReader}, \texttt{Arrays.asList()}, \texttt{HandlerAdapter} \\
Bridge & архитектура JDBC: интерфейсы и драйвер конкретной СУБД \\
Composite & граф сцены JavaFX; \texttt{Component} и \texttt{Container} в AWT \\
Decorator & \texttt{BufferedReader}, \texttt{GZIPOutputStream}, \texttt{Collections.unmodifiableList()} \\
Facade & \texttt{java.net.URL}; \texttt{JdbcTemplate}, \texttt{RestTemplate} \\
Flyweight & кеш \texttt{Integer.valueOf()}, пул строковых литералов \\
Proxy & \texttt{java.lang.reflect.Proxy}; \texttt{@Transactional}, \texttt{@Cacheable} \\
Chain of Responsibility & фильтры сервлетов; \texttt{SecurityFilterChain} \\
Command & \texttt{Runnable} и \texttt{Callable} в пуле потоков \\
Interpreter & \texttt{java.util.regex.Pattern}; язык выражений SpEL \\
Iterator & весь фреймворк коллекций; цикл \texttt{for-each} \\
Mediator & \texttt{ExecutorService}; \texttt{ApplicationEventPublisher} \\
Memento & сериализация объекта; точки сохранения внутри транзакции \\
Observer & слушатели JavaFX; \texttt{@EventListener} \\
State & \texttt{Thread.State}; проект Spring State Machine \\
Strategy & \texttt{Comparator}; \texttt{PasswordEncoder}, любой внедряемый бин \\
Template Method & \texttt{AbstractList}, \texttt{HttpServlet}; \texttt{JdbcTemplate} \\
Visitor & \texttt{FileVisitor}; альтернатива~--- \texttt{sealed} и \texttt{switch} \\
\end{longtable}

\section{Антипаттерны}

Антипаттерн отличается от обычной ошибки тем, что воспроизводится систематически
и разными людьми независимо друг от друга. Аналогия~--- беспорядок в квартире:
ни одна отдельно брошенная вещь не создаёт проблемы, но привычка «положу пока
сюда», повторённая триста раз, превращает квартиру в место, где ничего
невозможно найти.

\subsection{Божественный объект и выстрел дробью}

\term{God Object}~--- класс на полторы-три тысячи строк, в имени которого стоит
слово \texttt{Manager}, \texttt{Util} или \texttt{Processor}: двадцать полей и
полсотни методов, между которыми нет ничего общего, и все правки проекта так или
иначе задевают этот файл. Его невозможно понять целиком, протестировать по
частям и менять параллельно двум разработчикам. Лечится приёмом
\eng{Extract Class}: сгруппируйте методы по данным, с которыми они работают,
вынесите каждую группу в свой класс и передайте их друг другу через конструктор.

\term{Shotgun Surgery}~--- обратная беда: одно логическое изменение требует
мелких правок в десятке файлов. Добавили читателю поле «телефон»~--- правьте
сущность, объект передачи данных, преобразователь, форму, валидатор, три отчёта
и два теста; растёт вероятность что-то забыть, а забытое обнаруживается уже у
пользователя. Лечится сбором разбросанной логики в одном месте. Связь между
двумя антипаттернами прямая: в первом случае ответственность собрана слишком
густо, во втором размазана слишком тонко, а здоровое проектирование лежит
посередине~--- связанное вместе, несвязанное врозь.

\subsection{Спагетти-код и магические числа}

Признаки спагетти-кода~--- вложенность условий на пять-шесть уровней, методы по
двести строк и флаги вида \texttt{boolean isProcessed}, меняющие поведение
дальше по коду. Лечение~--- \term{ранние возвраты} (\eng{guard clauses}):
листинг~\ref{lst:11-13}, смысл тот же, вложенность нулевая.

\begin{listing}[H]
\begin{minted}{java}
// БЫЛО: логика утоплена в четырёх уровнях вложенности
public String checkOrder(Order order) {
    if (order != null) {
        if (order.getItems() != null) {
            if (!order.getItems().isEmpty()) {
                if (order.getTotal() > 0) {
                    return "ok";
                } else { return "нулевая сумма"; }
            } else { return "пустой заказ"; }
        } else { return "нет позиций"; }
    }
    return "нет заказа";
}

// СТАЛО: каждая проверка - отдельный выход
public String checkOrder(Order order) {
    if (order == null) return "нет заказа";
    if (order.getItems() == null) return "нет позиций";
    if (order.getItems().isEmpty()) return "пустой заказ";
    if (order.getTotal() <= 0) return "нулевая сумма";
    return "ok";
}
\end{minted}
\caption{Глубокая вложенность и ранние возвраты}
\label{lst:11-13}
\end{listing}

Практическое правило: если метод не помещается на экран, его нужно разбивать.

\term{Магические числа}~--- литералы без объяснения смысла. Через месяц ни автор,
ни коллега не скажут, что означает \texttt{3} в условии
\texttt{if (status == 3) return daysOverdue * 0.13 * 30;}, а при изменении
ставки придётся искать \texttt{0.13} по всему проекту и легко пропустить одно
вхождение. Лечение~--- перечисление \texttt{LoanStatus} вместо числового кода и
именованные константы \texttt{DAILY\_RATE} и \texttt{DAYS\_IN\_PERIOD}: ставка
меняется в одном месте, а невозможный статус отсекается компилятором.
Исключения есть: \texttt{0}, \texttt{1} и $-1$ в очевидном контексте
константами не оформляют.

\subsection{Ещё шесть ловушек}

\term{Copy-Paste Programming}~--- одинаковые блоки в пяти местах, отличающиеся
одной переменной: найденная ошибка исправляется в одном месте и остаётся в
четырёх. Лечится принципом \term{DRY} (\eng{Don't Repeat Yourself}): два
метода-близнеца, считающих среднее по разным спискам, сворачиваются в один
обобщённый \texttt{<T> double average(List<T> items, ToDoubleFunction<T> f)}.
Оговорка важна: DRY~--- про дублирование \emph{знания}, и куски кода, меняющиеся
по разным причинам, объединять не нужно. Рядом стоит \term{Feature Envy}~---
метод, который больше интересуется данными чужого класса, чем своими; узнаётся
по цепочкам геттеров вида \texttt{order.getClient().getAddress().getCity()},
из-за которых изменение структуры \texttt{Order} ломает чужой класс. Лечится
перемещением метода туда, где живут данные, и правилом \eng{«tell, don't ask»}.

\term{Lava Flow}~--- мёртвый код, который никто не решается удалить, от
закомментированных блоков «на всякий случай» до классов с пометкой «не трогать».
За него платят при каждом чтении и рефакторинге, ничего не получая взамен,
поэтому удаляйте: страховкой служит система контроля версий из
главы~\ref{ch:12}. \term{Golden Hammer}~--- «если единственный инструмент,
который у вас есть,~--- молоток, всё вокруг начинает выглядеть как гвоздь»:
команда, знающая только Hibernate, тянет ORM в ночную агрегацию пятидесяти
миллионов строк, где нужен обычный SQL; противоядие~--- явно сравнивать
альтернативы, как мы сравнивали JDBC и Hibernate в главе~\ref{ch:06}.

\term{Premature Optimization}~--- кеши и отказ от читаемых конструкций «ради
скорости» до того, как хоть что-то измерено: оптимизировать надо
\emph{измеренное} узкое место, а в веб-приложении это обычно проблема \eng{N+1}
запросов, когда ради ста заказов ORM делает ещё сто запросов за их клиентами.
Наконец, \term{Big Ball of Mud}~--- финальная стадия, к которой приводит
накопление всех предыдущих антипаттернов: архитектуры нет, слои перемешаны,
между пакетами циклические зависимости. Переписывание с нуля почти всегда
проваливается; работает другое: провести границы слоёв, покрыть тестами то, что
трогаете, и вытеснять старый код по кусочкам.

Пользуйтесь этим списком как чек-листом на код-ревью: сначала находите признак,
потом называете антипаттерн, и только потом предлагаете лечение.

\section{Рефакторинг: как применять паттерны на практике}

Квартиру не ремонтируют, снося стены наугад: сначала фотографируют, как было,
потом делают по одной комнате, и после каждой проверяют, что вода течёт, а свет
горит. Рефакторинг устроен так же: снимок текущего поведения, один маленький
шаг, проверка~--- и только потом следующий шаг. Посмотрим, как антипаттерны
лечатся паттернами на одном примере: метод
\texttt{calculate(String clientType, double amount)} состоит из цепочки
\texttt{else if} по строкам-числам, и при добавлении четвёртого типа клиента его
придётся править. Лечение идёт тремя шагами: убрать магические значения
перечислением, вынести каждый вариант расчёта в стратегию и связать их так,
чтобы новый тип клиента не требовал правки калькулятора
(листинг~\ref{lst:11-14}).

\begin{listing}[H]
\begin{minted}{java}
@FunctionalInterface
public interface PricingRule { double apply(double amount); }

public enum ClientType {
    REGULAR(amount -> amount),
    LOYAL(amount -> amount * 0.95),
    VIP(ClientType::vipPrice);

    private static final double VIP_THRESHOLD = 10_000;
    private final PricingRule rule;

    ClientType(PricingRule rule) { this.rule = rule; }

    // Правило вынесено в метод: статическое поле нельзя
    // читать прямо из объявления константы перечисления
    private static double vipPrice(double amount) {
        return amount > VIP_THRESHOLD
                ? amount * 0.85 : amount * 0.90;
    }

    public double priceFor(double amount) {
        return rule.apply(amount);
    }
}

public class PriceCalculator {
    // Не изменится, сколько бы типов клиентов ни добавили
    public double calculate(ClientType type, double amount) {
        return type.priceFor(amount);
    }
}
\end{minted}
\caption{Перечисление с поведением вместо цепочки условий}
\label{lst:11-14}
\end{listing}

Что мы получили: невозможный тип клиента отсекается компилятором, а не
исключением во время работы; каждое правило проверяется отдельным тестом;
добавление типа \texttt{PARTNER} требует одной строки в перечислении. Здесь
одновременно применены Стратегия и типичный для Java приём «перечисление с
поведением».

Порядок действий при рефакторинге складывается из четырёх правил. \emph{Тесты
первыми}: пока поведение не зафиксировано тестами из главы~\ref{ch:10}, любой
рефакторинг~--- риск. \emph{Мелкими шагами}: извлекли метод~--- запустили тесты,
переименовали~--- запустили. \emph{Один рефакторинг за раз}: не смешивайте
изменение поведения с изменением структуры в одном коммите, иначе при поиске
регрессии не поймёте, что сломалось. \emph{Пользуйтесь автоматикой среды
разработки}: извлечение метода и переименование выполняются за секунду и
безопасно. Практичный итог главы~--- таблица «что болит~--- что попробовать»
(табл.~\ref{tab:11-4}), читать которую нужно слева направо: сначала боль, потом
лекарство. И главное правило: паттерн вводят в ответ на реальную боль, а не
заранее~--- хороший код чаще получается рефакторингом простого решения, чем
проектированием сложного с нуля.

\begin{table}[htbp]
\caption{Выбор паттерна по симптому}
\label{tab:11-4}
\centering
\begin{tabular}{L{7.6cm}L{7.2cm}}
\toprule
Что болит & Что попробовать \\
\midrule
Длинный \texttt{if} или \texttt{switch} по типу либо статусу & Стратегия, Состояние, полиморфизм \\
Конструктор с восемью параметрами & Строитель \\
Класс зависит от конкретной реализации & DIP, внедрение зависимостей \\
Нужно добавить поведение, не меняя класс & Декоратор, Посетитель \\
Нужны логирование, кеш или права поверх вызова & Заместитель \\
Несовместимый сторонний интерфейс & Адаптер \\
Сложная подсистема, а нужен один сценарий & Фасад \\
«Кто-то должен узнать, что произошло событие» & Наблюдатель \\
Общий скелет алгоритма с разными шагами & Шаблонный метод \\
Древовидная структура, обрабатываемая единообразно & Компоновщик \\
Нужны отмена операции и история & Команда, Снимок \\
Много одинаковых объектов, не хватает памяти & Приспособленец \\
\bottomrule
\end{tabular}
\end{table}

\praktikum{}

Практикум устроен не так, как предыдущие: мы работаем с уже написанным кодом и
учимся видеть в нём структуру. Половина заданий начинается словами «вот код,
найдите в нём проблему», вторая половина~--- «а теперь исправьте»: врач сначала
ставит диагноз и только потом лечит. Всё занятие проходит на голой Java, без
систем сборки и фреймворков~--- так виднее сами паттерны. Создайте каталог
\texttt{patterns-lab} с подкаталогом \texttt{src}, а внутри~--- дерево пакетов
\texttt{ru/fa/patterns} с папками \texttt{solid}, \texttt{singleton},
\texttt{creational}, \texttt{behavioral} и \texttt{antipatterns}. Компилировать и
запускать будем командами из листинга~\ref{lst:11-15}: ключ
\texttt{-sourcepath} заставляет компилятор самому найти все нужные классы, а
\texttt{-encoding} обязателен~--- без него русские строки ломаются при
компиляции.

\begin{listing}[H]
\begin{minted}{bash}
javac -encoding UTF-8 -d out -sourcepath src \
      src/ru/fa/patterns/Setup.java
java -Dstdout.encoding=UTF-8 -cp out ru.fa.patterns.Setup
\end{minted}
\caption{Шаблон команд сборки и запуска}
\label{lst:11-15}
\end{listing}

\subsection*{Задание 1. SOLID: диагноз и лечение}

Наберите код из листинга~\ref{lst:11-16} в файл
\texttt{src/ru/fa/patterns/solid/ReaderService.java}; в этом файле ничего не
исправляйте~--- он нужен для отчёта. Понадобится ещё класс
\texttt{MySqlReaderStorage} с методом \texttt{insert(String name, String
email)}, печатающим имя и адрес.

\begin{listing}[H]
\begin{minted}{java}
public class ReaderService {

    private final List<String> readers = new ArrayList<>();
    private final MySqlReaderStorage storage =
            new MySqlReaderStorage();                  // метка 1

    public void register(String name, String email,
                         int categoryCode) {           // метка 2
        if (name == null || name.isBlank())
            throw new IllegalArgumentException("Пустое имя");
        readers.add(name);
        storage.insert(name, email);
        String body;                                   // метка 3
        if (categoryCode == 1) {
            body = "Обычный читатель, лимит книг 5.";
        } else if (categoryCode == 2) {
            body = "Студент, лимит книг 10.";
        } else {
            throw new IllegalArgumentException("Категория?");
        }
        System.out.println("SMTP -> " + email + ": " + body);
    }   // метка 4: отправка почты внутри регистрации

    public double fine(int categoryCode, int days) {   // метка 5
        if (categoryCode == 3) return 0;
        return days * 10.0;
    }

    // метка 6: сборка HTML-списка читателей
    public String toHtml() { /* ... StringBuilder ... */ }
}
\end{minted}
\caption{Код для разбора: шесть помеченных мест}
\label{lst:11-16}
\end{listing}

Заполните таблицу по строке на каждую метку: что здесь сделано, какой принцип
нарушен, чем это обернётся через полгода. Образец для метки~1: «сервис сам
создаёт конкретную реализацию хранилища~--- нарушен DIP~--- подменить хранилище
заглушкой в тесте невозможно».

Теперь лечение. В пакете \texttt{ru.fa.patterns.solid.refactored} заведите
запись \texttt{Reader} и перечисление \texttt{ReaderCategory} с константами
\texttt{REGULAR}, \texttt{STUDENT}, \texttt{TEACHER}, у каждой из которых свои
лимит книг, дневная пеня и метод \texttt{welcomeMessage()}; интерфейсы
\texttt{ReaderStorage} и \texttt{Notifier} с реализациями
\texttt{InMemoryReaderStorage} и \texttt{ConsoleNotifier};
\texttt{ReaderRegistrationService}, получающий обе зависимости через
конструктор; \texttt{FineCalculator} с методом
\texttt{fine(ReaderCategory c, int days)} без единого условия по категории;
\texttt{HtmlReaderList} и \texttt{SolidDemo}, который обходит
\texttt{ReaderCategory.values()} в цикле, регистрирует по читателю на категорию
и печатает пеню за четыре дня просрочки. Контрольное изменение: добавьте
четвёртую категорию~--- гостя с лимитом в две книги; правило приёмки~--- изменён
ровно один файл.

Третья часть~--- LSP и ISP. Создайте \texttt{LiskovDemo}, вложите в него
\texttt{Rectangle} и \texttt{Square} из листинга~\ref{lst:11-2} как статические
классы, добавьте метод \texttt{resizeAndCheck(Rectangle r)}, выставляющий
ширину~5 и высоту~4 и печатающий имя класса с площадью, и вызовите его для обоих
объектов: вторая строка разойдётся с ожиданием. Затем вылечите оба принципа~---
заведите \texttt{sealed interface Shape permits Rect, Sq} с методом
\texttt{int area()} и двумя реализациями-записями, а «толстый» интерфейс
устройства разбейте на \texttt{Printer}, \texttt{DocScanner} и
\texttt{FaxSender}.

Ответьте письменно: сколько у \texttt{ReaderService} причин для изменения; какая
площадь получилась у квадрата и почему клиентский код нельзя считать ошибочным;
почему неизменяемые объекты реже нарушают LSP.

\subsection*{Задание 2. Одиночка тремя способами}

Одиночка~--- единственный ключ от серверной: пока ключ один, порядок есть, но
если два сотрудника одновременно увидят, что ключа на доске нет, каждый закажет
себе новый. Ровно это происходит с наивным одиночкой в многопоточной среде.

В пакете \texttt{ru.fa.patterns.singleton} создайте класс \texttt{Counters} с
тремя статическими полями \texttt{NAIVE}, \texttt{SAFE} и \texttt{ENUM} типа
\texttt{AtomicInteger} и стенд \texttt{SingletonRace} с методом
\texttt{int distinctInstances(Supplier<Object> factory, int threads)}: он
запускает \texttt{threads} потоков, все они ждут общей отмашки на
\texttt{CountDownLatch}, одновременно вызывают \texttt{factory.get()} и
складывают объекты в множество на основе \texttt{IdentityHashMap}~--- оно
сравнивает по ссылке, а не по \texttt{equals}. Реализуйте три варианта:
\texttt{NaiveRegistry}~--- ленивый и небезопасный, с проверкой
\texttt{if (instance == null)} без синхронизации; \texttt{SafeRegistry}~---
двойная проверка блокировки из листинга~\ref{lst:11-3};
\texttt{EnumRegistry}~--- перечисление с единственной константой. В конструктор
каждого поставьте увеличение счётчика и \texttt{Thread.sleep(20)}: задержка
имитирует тяжёлую инициализацию и делает гонку заметной. Класс
\texttt{SingletonDemo} печатает для каждого варианта число разных объектов при
пятидесяти потоках и число вызовов конструктора.

Запустите демонстрацию пять раз подряд: у наивного варианта число объектов будет
меняться от запуска к запуску, но единицей не окажется почти никогда. Затем
уберите \texttt{Thread.sleep(20)} и запустите ещё пять раз~--- результат станет
непредсказуемым, иногда честно равным единице; это и есть самое неприятное
свойство гонок. Отдельно уберите слово \texttt{volatile} из
\texttt{SafeRegistry} и убедитесь, что тест по-прежнему показывает один объект.
Наконец, напишите \texttt{ReflectionAttack}: он получает конструктор
\texttt{NaiveRegistry} методом \texttt{getDeclaredConstructor()}, снимает защиту
вызовом \texttt{setAccessible(true)}, создаёт второй экземпляр, а затем пробует
проделать то же с \texttt{EnumRegistry}. Сведите три варианта в таблицу по
критериям: потокобезопасность, ленивость, устойчивость к рефлексии и к
сериализации, объём кода.

Ответьте письменно: в какой именно момент между проверкой и присваиванием второй
поток успевает вклиниться; какое исключение выдала виртуальная машина при
попытке продублировать перечисление; почему удаление \texttt{volatile} не
сломало программу и почему это не повод считать такой код правильным.

\subsection*{Задание 3. Строитель, фабричный и шаблонный методы}

Перенесите класс \texttt{Book} со Строителем из листинга~\ref{lst:11-4} в пакет
\texttt{ru.fa.patterns.creational} и доработайте: добавьте поля
\texttt{int pages} и \texttt{List<String> tags} (список в строителе объявите
сразу инициализированным, иначе книга без тегов уронит копирование с
\texttt{NullPointerException}), метод \texttt{tag(String)} и защитное
копирование \texttt{this.tags = List.copyOf(builder.tags)} в конструкторе.
В \texttt{build()} пустое название, пустой автор, год вне диапазона 1450--2100 и
неположительное число страниц должны давать \texttt{IllegalStateException} с
понятным сообщением. Класс \texttt{BuilderDemo} собирает три книги, печатает их,
а затем в \texttt{try/catch} пытается собрать книгу без автора и добавить
элемент в список тегов готовой книги, печатая пойманные исключения.

Вторая часть~--- экспорт каталога, где неизменны сбор текста и запись файла, а
переменен формат. Объявите интерфейс \texttt{CatalogFormatter} с методами
\texttt{String format(List<Book> books)} и \texttt{String fileExtension()} и
абстрактный класс \texttt{CatalogExporter} с абстрактным
\texttt{createFormatter()} и финальным \texttt{export}, который форматирует
книги, пишет результат через \texttt{Files.writeString} и печатает путь и
размер. Напишите две пары «форматировщик~--- экспортёр», для CSV и для Markdown,
и \texttt{ExporterDemo}, выгружающий четыре книги обоими. Здесь сразу два
паттерна: \texttt{createFormatter()}~--- Фабричный метод, а \texttt{export()}~---
Шаблонный, потому и объявлен \texttt{final}. Добавьте ещё
\texttt{CatalogExporters} со статическим методом \texttt{of(String format)},
возвращающим экспортёр через \texttt{switch}: это простая фабрика, не паттерн
GoF, но её постоянно путают с фабричным методом.

Ответьте письменно: что придётся создать и что изменить ради формата JSON~---
для фабричного метода и для простой фабрики; какой вариант соблюдает OCP; что
сломается, если подкласс переопределит \texttt{export}.

\subsection*{Задание 4. Стратегия, Наблюдатель и Декоратор}

Стратегия~--- это тарифы в такси: маршрут один, а считают по-разному. Объявите
функциональный интерфейс \texttt{FinePolicy} с методом
\texttt{double amount(int daysOverdue)} и перечисление \texttt{FineTariff},
хранящее по стратегии на тариф: \texttt{REGULAR}~--- десять рублей за день,
\texttt{STUDENT}~--- половина этой суммы, \texttt{TEACHER}~--- пеня всегда
нулевая, \texttt{VIP}~--- первые три дня бесплатно, дальше двойная ставка.
Правила выносите в приватные статические методы: читать статическое поле прямо
из объявления константы перечисления нельзя. Класс \texttt{StrategyDemo} печатает
пеню каждого тарифа за 0, 2, 5 и 30 дней просрочки и демонстрирует, что
\texttt{Comparator}~--- тоже стратегия: метод
\texttt{printSorted(List<Book> books, Comparator<Book> order)} вызовите трижды.

Вторая часть~--- Наблюдатель и его ловушки. Объявите функциональный интерфейс
\texttt{LoanListener} с методом \texttt{onIssued(String reader, String title)} и
класс \texttt{Library} со списком слушателей, подпиской, отпиской и методом
\texttt{issue}, который печатает сообщение о выдаче и обходит список. Порядок
подписки принципиален: первым идёт «разовый» слушатель, отписывающийся прямо
внутри обработчика (анонимный класс, иначе не сослаться на \texttt{this}),
вторым~--- слушатель, бросающий \texttt{RuntimeException}, третьим и
четвёртым~--- журнал выдач и счётчик в \texttt{AtomicInteger}. Вызовите
\texttt{issue} и почините две проблемы по очереди: сначала
\texttt{ConcurrentModificationException} от отписки внутри обхода (перебирайте
копию списка или храните слушателей в \texttt{CopyOnWriteArrayList}), затем
обрыв рассылки упавшим слушателем (\texttt{try/catch} вокруг каждого вызова).
После каждой починки запускайте программу и фиксируйте вывод.

Третья часть~--- Декоратор поверх стандартных потоков. Класс
\texttt{StreamChainDemo} тысячу раз записывает одну строку двумя способами: в
файл \texttt{catalog.txt.gz} через цепочку \texttt{BufferedWriter} $\to$
\texttt{OutputStreamWriter} $\to$ \texttt{GZIPOutputStream} $\to$
\texttt{FileOutputStream} и в \texttt{catalog.txt} без сжатия; напечатайте
размеры обоих файлов и прочитайте первую строку из архива зеркальной цепочкой.
Затем напишите декоратор \texttt{CountingOutputStream extends
FilterOutputStream}, считающий прошедшие байты (переопределить придётся и
\texttt{write(int)}, и \texttt{write(byte[], int, int)}), поставьте два таких
счётчика на разных уровнях матрёшки~--- до сжатия и после~--- и напечатайте
коэффициент сжатия, не забыв закрыть поток.

Ответьте письменно: почему исключение возникает при отписке первого слушателя, а
при отписке предпоследнего программа отработает молча и чем второй случай
опаснее; какой принцип SOLID стоит за тем, что новая реакция не требует правки
класса \texttt{Library}; какой класс в цепочке потоков Адаптер, а не Декоратор.

\subsection*{Задание 5. Антипаттерны: снимок поведения и рефакторинг по шагам}

Напишите в пакете \texttt{ru.fa.patterns.antipatterns} класс
\texttt{LibraryManager}, собравший в себе сразу несколько антипаттернов. Он
хранит книги и читателей в списках массивов \texttt{String[]} с комментариями о
том, что означает каждый элемент, а выданные книги~--- в
\texttt{Map<String, String>}. Метод \texttt{addBook} проверяет название, автора
и год четырьмя вложенными условиями с ветками \texttt{else} на каждом уровне;
\texttt{issue} печатает сообщение о выдаче и тут же отправляет письмо;
\texttt{fine} считает пеню по числовому коду категории с константами прямо в
коде; \texttt{averageBookYear} и \texttt{averageReaderCategory}~--- почти
дословные близнецы. Есть ещё \texttt{reportHtml}, \texttt{backup},
неиспользуемое поле \texttt{legacyCounter} и закомментированный
\texttt{exportToExcel}. Класс \texttt{LibraryManagerDemo} вызывает все методы с
фиксированными данными, включая заведомо неверные: пустое название, год 2200,
повторная выдача той же книги.

Найдите не менее шести антипаттернов и оформите таблицу «антипаттерн~--- где
именно~--- чем вреден~--- чем лечится». Отдельно разберите приём, не названный в
теории: доменные данные хранятся в \texttt{String[]} и \texttt{int} вместо
собственных типов, а компилятор не мешает перепутать первый элемент массива со
вторым; это одержимость примитивами, и лечат её записями. Затем снимите эталон
поведения: тестов у нас нет, поэтому их роль сыграет вывод программы,
перенаправленный в файл \texttt{before.txt} (флаг
\texttt{-Dstdout.encoding=UTF-8} обязателен: в файл Java пишет системной
кодировкой, и русские буквы превратятся в мусор). Рефакторинг ведите по шагам из
табл.~\ref{tab:11-5} в пакете \texttt{ru.fa.patterns.antipatterns.refactored},
после каждого шага пересобирая проект и сверяя вывод с эталоном.

\begin{table}[htbp]
\caption{Семь шагов рефакторинга}
\label{tab:11-5}
\centering
\begin{tabular}{cL{9.4cm}L{4cm}}
\toprule
Шаг & Что сделать & Против чего \\
\midrule
1 & переписать \texttt{addBook} на ранние возвраты & Spaghetti Code \\
2 & ввести константы \texttt{MIN\_YEAR}, \texttt{DAILY\_FINE} и перечисление \texttt{ReaderCategory} с явным кодом & Magic Numbers \\
3 & ввести записи \texttt{BookRecord} и \texttt{ReaderRecord}; массивы исчезают & одержимость примитивами \\
4 & свернуть методы-близнецы в обобщённый с \texttt{ToDoubleFunction} & Copy-Paste \\
5 & разложить класс на \texttt{BookCatalog}, \texttt{ReaderRegistry}, \texttt{LoanService}, \texttt{FineCalculator}, \texttt{HtmlCatalogReport} & God Object \\
6 & завести интерфейс \texttt{Notifier} и передавать его в конструкторе & нарушение DIP \\
7 & удалить \texttt{legacyCounter} и закомментированный метод & Lava Flow \\
\bottomrule
\end{tabular}
\end{table}

Числовой код категории на шаге~2 задайте явно, а не через \texttt{ordinal()}:
эталонная строка со средней категорией получена на кодах 1, 2 и 3, и опираться
здесь на порядок объявления констант~--- то самое сцепление с числом, от
которого шаг~2 как раз избавляет. Напишите \texttt{RefactoredDemo} с тем же
сценарием, сохраните его вывод в \texttt{after.txt} тем же способом и сравните
файлы: пустой результат сравнения означает, что рефакторинг прошёл чисто.

Ответьте письменно: сколько у \texttt{LibraryManager} ответственностей; какие
числа в коде магические; что придётся править, если завтра у книги появится
жанр; на каком шаге вывод разошёлся с эталоном и почему; строк кода стало
больше~--- почему это не аргумент против рефакторинга.

\subsection*{Результаты занятия}

По итогам занятия у вас должны быть готовы: таблица диагнозов по шести меткам
класса \texttt{ReaderService} и пакет с разложенными классами с подтверждением,
что новая категория добавилась правкой одного файла; вывод \texttt{LiskovDemo} и
его исправленная версия; три реализации Одиночки, результаты пяти запусков гонки
и сводная таблица вариантов; класс \texttt{Book} со Строителем и оба экспортёра
каталога; дополненное перечисление тарифов; все три версии класса
\texttt{Library}; разбор цепочки потоков и показания двух счётчиков; таблица
антипаттернов \texttt{LibraryManager}, файлы \texttt{before.txt} и
\texttt{after.txt} с результатом сравнения и журнал рефакторинга.

\kontrolvoprosy

\begin{enumerate}
\item Из каких четырёх элементов состоит описание паттерна и почему паттерн~---
      это не готовый код для копирования?
\item На какие три группы делятся паттерны «банды четырёх» и на какой вопрос
      отвечает каждая?
\item Что такое избыточное усложнение? Сформулируйте правило трёх и YAGNI.
\item Какая конструкция чаще всего сигнализирует о нарушении принципа
      открытости/закрытости?
\item Разберите на примере квадрата и прямоугольника, почему математическое
      «является» не совпадает с программным и почему ошибка здесь не в
      клиентском коде.
\item В чём разница между DIP, DI и IoC?
\item Зачем в двойной проверке блокировки нужно слово \texttt{volatile} и что
      может пойти не так без него?
\item Чем Фабричный метод отличается от Абстрактной фабрики и чем~--- от
      простой фабрики со \texttt{switch} с точки зрения OCP?
\item Какую проблему решает Строитель и почему собранный им объект обычно
      делают неизменяемым?
\item В цепочке \texttt{BufferedReader}, \texttt{InputStreamReader},
      \texttt{FileInputStream} какой класс Адаптер, какой Декоратор и по каким
      признакам вы их различили?
\item Чем Заместитель отличается от Декоратора при одинаковой технической
      структуре? Назовите четыре разновидности Заместителя.
\item Чем Стратегия отличается от Состояния и зачем шаблонный метод объявляют
      \texttt{final}?
\item Что такое God Object и Shotgun Surgery и почему их называют двумя
      сторонами одной ошибки?
\item Каков правильный порядок действий при рефакторинге и почему нельзя
      смешивать изменение структуры с изменением поведения в одном коммите?
\end{enumerate}

\testzadaniya

\begin{enumerate}

\item Что описано в книге «банды четырёх»?
  \begin{enumerate}
    \item[а)] 12 паттернов, разделённых на классовые и объектные
    \item[б)] 23 паттерна, разделённых на порождающие и поведенческие
    \item[в)] 23 паттерна, разделённых на порождающие, структурные и
              поведенческие
    \item[г)] 5 принципов SOLID и 18 паттернов уровня архитектуры
  \end{enumerate}

\item Метод \texttt{area(Object shape)} состоит из длинной цепочки
      \texttt{if (shape instanceof ...)}. Симптом нарушения какого принципа?
  \begin{enumerate}
    \item[а)] SRP~--- класс делает слишком много
    \item[б)] ISP~--- интерфейс навязывает лишние методы
    \item[в)] LSP~--- подкласс не подставим вместо родителя
    \item[г)] OCP~--- класс не закрыт для изменения
  \end{enumerate}

\item Для классов из листинга~\ref{lst:11-2} выполнили
      \texttt{Rectangle r = new Square(); r.setWidth(5); r.setHeight(4);}.
      Что вернёт \texttt{r.area()}?
  \begin{enumerate}
    \item[а)] 20
    \item[б)] 16
    \item[в)] 25
    \item[г)] код не скомпилируется
  \end{enumerate}

\item В реализации интерфейса три метода из пяти бросают
      \texttt{UnsupportedOperationException}. Нарушение какого принципа это
      выдаёт?
  \begin{enumerate}
    \item[а)] ISP~--- принцип разделения интерфейсов
    \item[б)] SRP~--- принцип единственной ответственности
    \item[в)] DIP~--- принцип инверсии зависимостей
    \item[г)] OCP~--- принцип открытости/закрытости
  \end{enumerate}

\item Зачем поле \texttt{instance} в Одиночке с двойной проверкой блокировки
      объявляют \texttt{volatile}?
  \begin{enumerate}
    \item[а)] чтобы поле стало доступно из других пакетов
    \item[б)] чтобы запретить переупорядочение операций: иначе другой поток
              может увидеть ссылку на ещё не достроенный объект
    \item[в)] чтобы поле хранилось в стеке потока, а не в куче
    \item[г)] чтобы синхронизировать доступ и обойтись без блока
              \texttt{synchronized}
  \end{enumerate}

\item Какой метод стандартной библиотеки является настоящим Фабричным методом
      GoF?
  \begin{enumerate}
    \item[а)] \texttt{Calendar.getInstance()}
    \item[б)] \texttt{Collection.iterator()}
    \item[в)] \texttt{Integer.valueOf()}
    \item[г)] \texttt{HttpRequest.newBuilder()}
  \end{enumerate}

\item Что вернёт \texttt{new BracketsDecorator(new UpperCaseDecorator(new
      FileDataSource())).read()}, если \texttt{FileDataSource.read()}
      возвращает «данные отчёта»?
  \begin{enumerate}
    \item[а)] данные отчёта
    \item[б)] [ДАННЫЕ ОТЧЁТА]
    \item[в)] [данные отчёта]
    \item[г)] ДАННЫЕ ОТЧЁТА
  \end{enumerate}

\item Почему для \texttt{Integer a = 127, b = 127} сравнение \texttt{a == b}
      даёт \texttt{true}, а для 128~--- \texttt{false}?
  \begin{enumerate}
    \item[а)] \texttt{Integer.valueOf()} кеширует значения от $-128$ до
              $127$~--- это Приспособленец
    \item[б)] значения до 127 умещаются в байт и сравниваются по значению
    \item[в)] автоупаковка для чисел меньше 128 не выполняется
    \item[г)] виртуальная машина оптимизирует \texttt{==} для малых чисел
  \end{enumerate}

\item Чем Состояние отличается от Стратегии, если структурно они почти
      одинаковы?
  \begin{enumerate}
    \item[а)] Стратегия~--- структурный паттерн, Состояние~--- поведенческий
    \item[б)] Стратегия построена на наследовании, Состояние~--- на композиции
    \item[в)] Состояние применимо лишь к перечислениям
    \item[г)] в Стратегии алгоритм задаёт клиент извне, а в Состоянии
              следующий переход выбирает сам объект состояния
  \end{enumerate}

\item Какую роль играет \texttt{InputStreamReader} в цепочке
      \texttt{BufferedReader}~--- \texttt{InputStreamReader}~---
      \texttt{FileInputStream}?
  \begin{enumerate}
    \item[а)] Адаптера: превращает байтовый поток в символьный
    \item[б)] Декоратора: добавляет метод \texttt{readLine()}
    \item[в)] Фасада: прячет работу с файловой системой
    \item[г)] Заместителя: контролирует доступ к файлу
  \end{enumerate}

\item Добавление одного поля требует правок в сущности, объекте передачи
      данных, преобразователе, форме, валидаторе и трёх отчётах. Что это?
  \begin{enumerate}
    \item[а)] God Object, а Shotgun Surgery~--- его частный случай
    \item[б)] Copy-Paste Programming, лечится обобщениями
    \item[в)] Big Ball of Mud, лечится переписыванием с нуля
    \item[г)] Shotgun Surgery~--- обратная сторона God Object:
              ответственность размазана слишком тонко
  \end{enumerate}

\end{enumerate}
